\documentclass[10pt,a4paper]{amsart}
\usepackage[margin=19mm]{geometry}
\usepackage{amsmath,amssymb,mathtools}
\usepackage[scr=rsfs,cal=euler]{mathalfa}
\usepackage{xfrac}
\usepackage[unicode,hidelinks,bookmarks=false]{hyperref}
\newcommand{\ie}{i.e.\ }
\newcommand{\cf}{cf.\ }
\newcommand{\resp}{respectively\ }
\newcommand{\Subsection}{\subsection}
\providecommand{\nobreakdash}{\nobreak\hbox{-}}
\setlength{\emergencystretch}{2em}
\multlinegap0pt
\usepackage{xcolor}
\usepackage{amssymb}
\usepackage{mathtools}
\usepackage[shortlabels]{enumitem}
\usepackage{bm}
\usepackage{tikz}
\newtheorem{theorem}{Theorem}
\title{On the Architectural Complexity of Neural Networks}
\author{Nicholas J. Cooper, Fran\c{c}ois G. Meyer, Michael L. Roberts, Carlos Zapata-Carratal\'a, Lijun Chen, Danna Gurari}
\date{Source: arXiv:2605.04325v1 (CC BY 4.0)}
\begin{document}
\maketitle

\noindent\emph{Source and licence.} On the Architectural Complexity of Neural Networks. Version arXiv:2605.04325v1 (CC BY 4.0), distributed by arXiv under the Creative Commons Attribution 4.0 International licence, http://creativecommons.org/licenses/by/4.0/, as recorded in the arXiv licence metadata for this exact version and shown on the abstract page https://arxiv.org/abs/2605.04325v1. Full licence notice: \texttt{licenses/arxiv-2605.04325-CC-BY-4.0.txt}.\par\smallskip

\noindent\textit{Editorial scope.} Counted unit: the article's theorem thrm:CanonicalRepsExist (source0.tex lines 1051--1054). The article's complete original proof of it is delivered with it (source0.tex lines 1055--1089, one \texttt{proof} environment of 35 lines). No mathematics is written, repaired or re-derived: the statement and the proof are the article's own lines, in the article's own notation, and no retained line is substituted or re-flowed.  No further source-local context is needed: every cross-reference of every retained line resolves inside the two delivered spans.  Nothing was omitted from any retained span, so the package carries no omission record.  No retained line contains a citation, so the wrapper carries no bibliography and no bibliography entry.  No retained line contains a figure environment, an image-inclusion command or drawing code, so no image is delivered and the package declares no asset dependency.  Editorial formatting only, in addition: an AMS \texttt{amsart} preamble that restates the article's own declarations and macros used by the retained lines, namely the environments \texttt{theorem} on separate counters, together with the standard packages those lines need.  The retained environments print in this document as Theorem 1 in order of appearance, whereas the article prints the same items with the numbering of its own full document.  The complete native source of the article as distributed in the arXiv e-print for this exact version is bundled unchanged as \texttt{sources/199849/source0.tex}.
\begin{theorem}
	Let $s^3$ be a generalized tensor which satisfies the strong SOCCs. There exists a procedure to produce a canonical representative $s^3_c$ for $s^3$. Furthermore, $s^3_c$ is unique up to permutations of its $2$-cells.
	\label{thrm:CanonicalRepsExist}
\end{theorem}

\begin{proof}
	Let $s^3$ be a generalized tensor, and let $s^3_m$ be a maximal representative for $s^3$. We will construct a canonical representative of $s^3$ by combining overlapping $1$-cells from each $s^2 \in s^3_m$. That is, we are interested in cases when:
	\begin{equation*}
		\begin{aligned}
			s^1_i, s^1_j \in s^2 \in s^3_m \text{ satisfy: }& \big( \nexists s^1_k \text{ with } s^1_i \subsetneq s^1_k \text{ or } s^1_j \subsetneq s^1_k \big) \ \wedge \ \big( s^1_i \cap s^1_j \neq \emptyset \big)\\
		\end{aligned}
	\end{equation*}
	We recall that such $\subsetneq$-maximal $1$-cells correspond to overlapping hyper-edges of the associated PWO-HG.
	Claim: if there are no such $s^1_i, s^1_j \in s^2 \in s^3_m$, then $s^3_m$ is already a canonical representative for $s^3$.\\
	We will show that no $1$-cell can be removed from $s^3_m$. To see this, recall that each $0$-cell of a generalized tensor must be contained in the intersection of a transversal of the partition of $\subsetneq$-maximal $1$-cells. In particular, this means that the claim's hypothesis implies that each $2$-cell must itself be a partition of the $0$-cells. It is then clear that no $1$-cell can removed from $s^3_m$, as otherwise some $2$-cell would no longer cover the $0$-cells. This proves the claim.
	
	We now assume that there exist some overlapping $\subsetneq$-maximal $s^1_i, s^1_j \in s^2 \in s^3_m$ and consequently we must produce a procedure for merging the cells. So, fix $b \in s^1_i \cap s^1_j$ and let $index_i(\cdot): s^1_i \rightarrow \big \langle \{1, 2, ..., |s^1_i|\}, < \big \rangle$ be the index function for $s^1_i$. We extend $index_i(\cdot)$ to a new function $index()$ defined on $s^1_i \cup s^1_j$ by the following rule:
	\begin{equation*}
		index(a) = \begin{cases}
			index_i(a), \ a \in s^1_i\\
			index_i(b) + d_{s^1_j}(\gamma_{b, a}), \text{ otherwise}
		\end{cases}
	\end{equation*}
	Claim: $index(\cdot)$ is a well-defined function on $s^1$.\\
	Suppose not. The only way $index(\cdot)$ could assign two different values to some $a \in s^1$ is if $a \in s^1_i \cap s^1_j$ and the two cases of the definition disagree. This means that $index_i(a) \neq index_i(b) + d_{s^1_j}(\gamma_{b, a})$. So, $index_i(a) - index_i(b) \neq d_{s^1_j}(\gamma_{b, a})$. But, $index_i(a) - index_i(b) = d_{s^1_i}(\gamma_{b, a})$, meaning that $a$ and $b$ have paths of different distances within the same $2$-cell. This contradicts the strong SOCC, proving the claim.
	
	Next, define a relation on $s^1 := s^1_i \cup s^1_j$ based on $index(\cdot)$, that is:
	\begin{equation*}
		\forall a, c \in s^1, a <_{s^1} c \iff index(a) < index(c)
	\end{equation*}
	Claim: $<_{s^1}$ is a strict weak order.\\
	By virtue of being defined by a function to the integers with $<$, $<_{s^1}$ is automatically transitive and has a transitive incomparability relation. We must now verify that $<_{s^1}$ is irreflexive and asymmetric. So, assume for the sake of contradiction that $<_{s^1}$ is not asymmetric. Then, we can find some $a, b \in s^1$ with $a <_{s^1} b$ and $b <_{s^1} a$. By construction, this means that $index(a) < index(b)$ and $index(b) < index(a)$. This is impossible because $index(\cdot)$ is well-defined. An analogous argument shows that $<_{s^1}$ is irreflexive. This proves the claim.
	
	Claim: Replacing $s^1_i$ and $s^1_j$ with $s^1$ does not change any path distances on the PWO-HG associated to $s^3$.\\
	Let $\gamma_{a, c}$ be a path between $a, c \in s^1$. By construction, if $a, c \in s^1_i$, then $d_{s^1}(\gamma_{a, c}) = d_{s^1_i}(\gamma_{a, c})$. Similarly, if $a, c \in s^1_j$, $d_{s^1}(\gamma_{a, c}) = d_{s^1_j}(\gamma_{a, c})$. So, we now assume that $a \in s^1_i$ and $c \in s^1_j$. Because $b \in s^1_i \cap s^1_j$, we can decompose $\gamma_{a, c}$ into two paths $\gamma_{a, c} = \gamma_{b, c} \circ \gamma_{a, b}$. Because $a, b \in s^1_i$, $d_{s^1}(\gamma_{a, b}) = d_{s^1_i}(\gamma_{a, b})$. Because $b, c \in s^1_j$, $d_{s^1}(\gamma_{b, c}) = d_{s^1_j}(\gamma_{b, c})$. So, $d_{s^1}(\gamma_{a, c}) = d_{s^1_i}(\gamma_{a, b}) + d_{s^1_j}(\gamma_{b, c})$. We conclude that the generalized tensor formed from $s^3_m$ by deleting $s^1_i$ and $s^1_j$ and then adding $s^1$ is equivalent to $s^3_m$. This shows that this process produces a canonical representative $s^3_c$ for $s^3$. This proves the claim.
	
	We must now show that $s^3_c$ is unique up to mode permutations. So, let $s^3_{c'}$ be a distinct canonical representative for $s^3_m$. Because equivalence of generalized tensors is transitive, $s^3_c$ and $s^3_{c'}$ are equivalent so they cannot differ in their $0$-cells. Let $g$ be the map between the $2$-cells of $s^3_c$ and $s^3_{c'}$. By a previous argument, the $2$-cells of each representative must partition the $0$-cells. As we need only prove uniqueness up to permutations of $2$-cells, there must exist some $s^2_c \in s^3_c$ such that $s^2_{c'} = g(s^2_c) \neq s^2_c$. In order for $s^2_c$ and $s^2_{c'}$ to be distinct partitions of the same set, there must be $0$-cells $a, b$ which belongs to the same $s^1 \in s^2_c$ and such that $a \in s^1_1 \in s^2_{c'}$ and $b \in s^1_2 \in s^2_{c'}$. By connectivity, let $\gamma_{a, b}$ be a path between $a$ and $b$ in $s^3_{c'}$. There are two possibilities: either $d_{s^2}(\gamma_{a, b}) \neq 0$ for some $s^2 \neq s^2_{c'}$, or $d_{s^2}(\gamma_{a, b}) = 0$ $\forall s^2 \neq s^2_{c'}$. The first case is impossible because there is a direct path between $a$ and $b$ in $s^3_c$ which only passes through $s^2_c$, so this would violate the fact that $s^3_c$ and $s^3_{c'}$ are equivalent. The second case decomposes into two more cases. If $|d_{s^2_{c'}}(\gamma_{a, b})| \leq 1$, then by the maximiality of $s^3_{c'}$, $a, b$ would belong to the same $1$-cell in $s^2_{c'}$ which is a contradiction. So, we now assume that $|d_{s^2_{c'}}(\gamma_{a, b})| > 1$ and $d_{s^2}(\gamma_{a, b}) = 0$ for all $s^2 \neq s^2_{c'}$. Because $s^3_c$ and $s^3_{c'}$ are equivalent, this means that there must exist some $0$-cells $x_1, x_2, ..., x_n$ such that $[a]_{s^1} <_{s^1} [x_1]_{s^1} <_{s^1} [x_2]_{s^1} ... <_{s^1} [x_n]_{s^1} <_{s^1} [b]_{s^1}$. By the maximality of $s^3_{c'}$, this means that $x_1 \in s^1_1$ and $x_n \in s^1_2$. Analogously, $x_2 \in s^1_1$ and $x_{n-1} \in s^1_2$, and so on. We can continue this process until we reach some $i$ such that $d_{s^1}(\gamma_{x_i, x_{i+1}}) = 1$ and $x_i \in s^1_1$ and $x_{i+1} \in s^1_2$. By the previous argument, this is impossible.
	
	We conclude that $s^3_c$ is unique up to mode permutations.
\end{proof}

\end{document}
